\documentclass[11pt]{amsart}
\usepackage[margin=1in]{geometry}
\usepackage{amsmath,amssymb,amsthm,mathtools}
\usepackage[colorlinks=true,linkcolor=blue,citecolor=blue,urlcolor=blue]{hyperref}
\usepackage{xurl}
\usepackage{booktabs}
\usepackage{enumitem}

\newcommand{\R}{\mathbb{R}}
\newcommand{\C}{\mathbb{C}}
\newcommand{\Z}{\mathbb{Z}}
\newcommand{\Lam}{\Lambda_M}

\theoremstyle{plain}
\newtheorem{theorem}{Theorem}[section]
\newtheorem{corollary}[theorem]{Corollary}
\newtheorem{proposition}[theorem]{Proposition}
\newtheorem{lemma}[theorem]{Lemma}
\theoremstyle{definition}
\newtheorem{definition}[theorem]{Definition}
\theoremstyle{remark}
\newtheorem{remark}[theorem]{Remark}
\newtheorem{openproblem}[theorem]{Open Problem}

\title[Exact triad structure of truncated 3-D Navier--Stokes]
{Machine-checked triad identities for the Fourier--Galerkin truncation\\
of the 3-D Navier--Stokes equations, with the vanishing locus\\
of the helical interaction coefficient in closed form}

\author[Preprint]{Author list and affiliations to be supplied}
\date{\today}

\begin{document}
\maketitle

\begin{abstract}
We give a formally verified account of the quadratic nonlinearity of the incompressible 3-D
Navier--Stokes equations restricted to the Fourier--Galerkin truncation
$\Lam = \{k \in \Z^3 : 0 < |k| \le M\}$. (i)~A single \emph{weighted triad identity} yields
exact energy conservation (constant weight) and the exact enstrophy production
(weight $|k|^2$) as two instances of one computation, isolating the production mechanism in the
factor $|r|^2 - |q|^2$, the wavenumber disparity of a triad's two swapped members. (ii)~For
Waleffe's helical decomposition we give the triad interaction coefficient in a square-root-free
integer frame, prove that the resonance condition $s_k|k| + s_p|p| + s_q|q| = 0$ holds for some
signs if and only if the triad is collinear (a fact of Euclidean geometry, independent of the
lattice), and prove the ``if and only if'': the coefficient vanishes exactly on the collinear
triads and on the balance locus $s_p|p| = s_q|q|$, and nowhere else. (iii)~On any periodic lattice
without $2$-torsion the resonant-triad incidence structure has an exactly solvable spectrum with
normalised gap $\to 1$; on the spherical truncation the gap appears to converge to $5/6$, which we
state as an open conjecture with six exact data points. Every theorem is machine-checked in Lean~4
against Mathlib with axiom footprint $\{\mathtt{propext}, \mathtt{Classical.choice},
\mathtt{Quot.sound}\}$ and no \texttt{sorry}; every closed form is independently cross-checked in
exact integer arithmetic. The identities and the coefficient are classical (Kraichnan, Waleffe);
what is new is the unified, machine-checked form, the explicit converse, and the exact spectrum.
All results are at fixed $M$; whether the production admits any bound uniform in $M$ is not
addressed.
\end{abstract}

\section{Introduction}

The Fourier--Galerkin truncation of the incompressible 3-D Navier--Stokes equations retains the
Fourier coefficients $u_k$, $k \in \Z^3$, with $0 < |k| \le M$ and discards the rest. The result
is a finite-dimensional system whose global well-posedness is elementary, and whose quadratic
nonlinearity keeps, exactly, the triad structure $p + q = k$ that carries the energy cascade and
vortex stretching in the full equations. Detailed conservation over triads goes back to
Kraichnan~\cite{kraichnan1959}; the helical decomposition, in which each mode splits into two
chiralities and each triad interaction reduces to one complex coefficient, is due to
Waleffe~\cite{waleffe1992}. Nothing in this paper changes those facts. What it adds is a
\emph{formally verified} version of them, organised so that the structure is visible in the
proofs rather than in the commentary, together with two statements we have not seen made
explicit: an ``if and only if'' for the vanishing of the helical coefficient
(Theorem~\ref{thm:converse}), and the exact spectrum of the resonant-triad incidence structure on
a periodic lattice (Theorem~\ref{thm:torus}), whose spherical analogue we leave as a sharply
posed conjecture.

The organising observation of Section~\ref{sec:weighted} is that energy conservation and enstrophy
production are the constant-weight and $|k|^2$-weight instances of one identity on triads
(Theorem~\ref{thm:weighted}): everything common to a triad's two orderings cancels, and what
remains is the difference of the weight at the two swapped members. For the energy weight the
difference is zero, which is the conservation law; for the enstrophy weight it is the wavenumber
disparity $|r|^2 - |q|^2$, which is the production mechanism, exactly.

Sections~\ref{sec:helical}--\ref{sec:closed} treat the helical coefficient. Working in a frame
built from the integer vector $p \times q$, and with helical vectors left unnormalised, the
coefficient becomes a polynomial identity in integer variables and three formal magnitudes, so
that it can be checked in exact integers over thousands of lattice triads --- which caught a sign
error in our first derivation that a formal proof would have certified just as happily
(Remark~\ref{rem:sign}). Section~\ref{sec:torus} gives the spectral result and the conjecture;
Section~\ref{sec:open} collects open questions.

\paragraph{Verification.} Every theorem stated below without qualification is proved in the
Lean~4 proof assistant~\cite{lean4} against a pinned revision of Mathlib~\cite{mathlib}. For each,
the compiled theorem's axiom footprint, reported by \texttt{\#print axioms}, is exactly
\{\texttt{propext}, \texttt{Classical.choice}, \texttt{Quot.sound}\} --- the standard axioms of
Mathlib itself --- and there is no use of \texttt{sorry}. Closed forms involving the irrational
magnitudes $|k|$ are additionally checked by an independent exact-integer computation
(Remark~\ref{rem:sign}). The development is public (Appendix).

\section{Setup}
\label{sec:setup}

Fix $M \ge 1$ and let $\Lam := \{k \in \Z^3 : 0 < |k|^2 \le M^2\}$, closed under $k \mapsto -k$.
A \emph{state} is $u = (u_k)_{k \in \Z^3}$, $u_k \in \C^3$, with $u_k = 0$ for $k \notin \Lam$,
$u_0 = 0$, and
\[
  \text{(divergence-free)}\quad k \cdot u_k = 0, \qquad
  \text{(real)}\quad u_{-k} = \overline{u_k}, \qquad \text{for all } k.
\]
Here $a \cdot b = \sum_i a_i b_i$ is the \emph{bilinear} dot product (no conjugation) and
$\langle a, b\rangle = \sum_i \overline{a_i}\, b_i$ the Hermitian pairing. Let $\mathbb{P}(k)$ be
the orthogonal projector onto $k^\perp$. The truncated Navier--Stokes nonlinearity is
\begin{equation}
  \label{eq:Bop}
  B(u,v)_k \;=\; \mathbb{P}(k)\Bigl[-i\!\!\sum_{\substack{p+q=k\\ p,q \in \Lam}}\!\!(q\cdot u_p)\,v_q\Bigr],
  \qquad k \in \Lam,
\end{equation}
and the truncated viscous system is $\dot u_k = -\nu|k|^2 u_k + B(u,u)_k$, $\nu \ge 0$. Energy
and enstrophy are $E(u) = \sum_k |u_k|^2$ and $\Omega(u) = \sum_k |k|^2|u_k|^2$.

\emph{Triad sums.} Let $T_M := \{(p,q) \in \Lam \times \Lam : -(p+q) \in \Lam\}$, and for
$(p,q) \in T_M$ write $r := -(p+q)$, so $p + q + r = 0$. All ``sums over triads'' below are sums
over the ordered pairs $(p,q) \in T_M$; each unordered triad is therefore counted once per ordered
choice of its first two members. The swap $\sigma(p,q) := (p, r)$ is an involution of $T_M$ that
fixes $p$ and exchanges $q$ and $r$.

\section{The weighted triad identity, and its two instances}
\label{sec:weighted}

\begin{theorem}[Weighted triad identity]
\label{thm:weighted}
Let $K$ be a field in which $2 \ne 0$, let $k_p, k_q, k_r, u_p, u_q, u_r \in K^3$ with
$k_p + k_q + k_r = 0$ and $k_p \cdot u_p = 0$, and let $w_q, w_r \in K$. Then
\begin{equation}
  \label{eq:weighted}
  w_r\,(k_q \cdot u_p)(u_q \cdot u_r) \;+\; w_q\,(k_r \cdot u_p)(u_r \cdot u_q)
  \;=\; (w_r - w_q)\,(k_q \cdot u_p)(u_q \cdot u_r).
\end{equation}
Consequently, for a state $u$ on $\Lam$ (divergence-free suffices; reality is not used) and any
weight $w : \Lam \to \C$,
\begin{equation}
  \label{eq:wsum}
  2\sum_{(p,q) \in T_M} w(r)\,(q \cdot u_p)(u_q \cdot u_r)
  \;=\; \sum_{(p,q) \in T_M} \bigl(w(r) - w(q)\bigr)(q \cdot u_p)(u_q \cdot u_r).
\end{equation}
\end{theorem}

\begin{proof}
From $k_p + k_q + k_r = 0$ and $k_p \cdot u_p = 0$, $k_r \cdot u_p = -(k_q \cdot u_p)$; the dot
product is symmetric, so $u_r \cdot u_q = u_q \cdot u_r$. Substituting gives~\eqref{eq:weighted}.
For~\eqref{eq:wsum}, the involution $\sigma$ preserves $T_M$, so the left side is the sum over
$T_M$ of the summand plus its $\sigma$-image, which is~\eqref{eq:weighted} term by term with
$k_p = p$, $k_q = q$, $k_r = r$ (divergence-freeness at $p$ is exactly $p \cdot u_p = 0$).
\end{proof}

Divergence-freeness at $p$ is used exactly once and is necessary: a non-solenoidal field breaks
the identity (this is one of the formal development's negative controls).

\begin{corollary}[Exact energy conservation]
\label{cor:energy}
For every real, divergence-free state $u$ on $\Lam$, and every $M$,
\begin{equation}
  \label{eq:energy}
  \operatorname{Re} \sum_{k \in \Lam} \langle u_k, B(u,u)_k\rangle \;=\; 0 .
\end{equation}
\end{corollary}

\begin{proof}
Since $u_k \perp k$, the projector drops out of the outer pairing, and reality
$\overline{u_k} = u_{-k}$ turns the Hermitian pairing into the bilinear one:
$\sum_k \langle u_k, B_k\rangle = -i \sum_{(p,q) \in T_M} (q \cdot u_p)(u_q \cdot u_r)$ with
$r = -(p+q) = -k$. This is~\eqref{eq:wsum} with $w \equiv 1$, whose right side vanishes; as
$2 \ne 0$ in $\C$, the sum is $0$.
\end{proof}

\begin{corollary}[Exact enstrophy production]
\label{cor:enstrophy}
For every real, divergence-free state $u$ on $\Lam$,
\begin{equation}
  \label{eq:enstrophy}
  2 \sum_{k \in \Lam} |k|^2 \langle u_k, B(u,u)_k \rangle
  \;=\; -i \sum_{(p,q) \in T_M} \bigl(|r|^2 - |q|^2\bigr)\,(q \cdot u_p)(u_q \cdot u_r),
  \qquad r = -(p+q).
\end{equation}
The physical production rate $\frac{1}{2}\frac{d}{dt}\Omega\big|_{\text{nonlinear}}$ is the real
part of the left side.
\end{corollary}

\begin{proof}
As in Corollary~\ref{cor:energy}, with the weight $w(r) = |r|^2 = |k|^2$ attached to the third
member; apply~\eqref{eq:wsum}.
\end{proof}

Equations~\eqref{eq:energy} and~\eqref{eq:enstrophy} are the same computation with and without the
weight. The whole difference between conservation and production is the factor $|r|^2 - |q|^2$: a
triad transfers enstrophy in proportion to the disparity in wavenumber of its two swapped members,
and none if they share a sphere. The formal development also checks that the right side
of~\eqref{eq:enstrophy} is nonzero on an explicit Galerkin state, so the identity is not vacuous.

\begin{corollary}[Balance laws]
\label{cor:balance}
Let $F_k = -\nu|k|^2 u_k + B(u,u)_k$ with $\nu \ge 0$. For every real, divergence-free state,
\begin{align}
  \operatorname{Re} \sum_k \langle u_k, F_k \rangle &= -\nu \sum_k |k|^2 |u_k|^2 \;\le\; 0,
  \label{eq:energybalance}\\
  \operatorname{Re} \sum_k |k|^2 \langle u_k, F_k \rangle
  &= -\nu \sum_k |k|^4 |u_k|^2 \;+\; \operatorname{Re} \sum_k |k|^2\langle u_k, B(u,u)_k \rangle.
  \label{eq:enstrophybalance}
\end{align}
\end{corollary}

Equation~\eqref{eq:energybalance} says the truncated energy is non-increasing;
\eqref{eq:enstrophybalance} exhibits the production term of~\eqref{eq:enstrophy} as the sole
obstruction to the same statement for the enstrophy. Nothing here bounds that term against the
dissipation $\nu \sum_k |k|^4 |u_k|^2$, at fixed $M$ or uniformly; see Section~\ref{sec:open}.

\section{The helical decomposition and its interaction coefficient}
\label{sec:helical}

Following Waleffe~\cite{waleffe1992}, each transverse plane $k^\perp$ is spanned by two
helical vectors $h^{\pm}(k)$ with $i\,k \times h^{s}(k) = s|k|\,h^{s}(k)$, $s \in \{+1,-1\}$. We
fix them in a way suited to exact arithmetic.

\begin{definition}[Helical vectors in a frame]
\label{def:h}
For an integer vector $N \perp k$, $N \ne 0$, set
\[
  h^{s}_N(k) \;:=\; N \times k \;+\; i\,s\,|k|\,N \;\in\; \C^3, \qquad s \in \{\pm 1\}.
\]
\end{definition}

For $N \perp k$ these satisfy $h^s_N(k)\cdot h^s_N(k) = 0$, $\overline{h^s_N(k)} = h^{-s}_N(k)$,
$k \cdot h^s_N(k) = 0$ and $i\,k \times h^s_N(k) = s|k|\,h^s_N(k)$ (all machine-checked). They are
Waleffe's $h^s(k)$ multiplied by $|N||k|$: keeping them unnormalised removes every square root and
every division from the proofs (in a proof assistant, $x/0 = 0$ silently proves a different
theorem). For a triad $p + q = k$ we use the \emph{common} frame vector $N := p \times q$, which is
an integer vector orthogonal to $p$, $q$, and $k$, and vanishes exactly on collinear triads.

Expanding $u_k = \sum_s a^s_k h^s(k)$, projecting $B(u,u)_k$ onto $h^{s_k}(k)$, and antisymmetrising
the two orderings of the triad (the same swap as in Theorem~\ref{thm:weighted}) produces, for each
triad and each chirality triple $(s_k, s_p, s_q)$, the coefficient
\begin{equation}
  \label{eq:Cdef}
  C \;=\; -\tfrac14\,\bigl(s_p|p| - s_q|q|\bigr)\, g, \qquad
  g \;=\; \bigl(h^{s_p}_N(p) \times h^{s_q}_N(q)\bigr)\cdot \overline{h^{s_k}_N(k)},
\end{equation}
so that the helical amplitude equations read $\dot a^{s_k}_k = \sum_{p+q=k}\sum_{s_p,s_q} C\,
a^{s_p}_p a^{s_q}_q + \dots$ up to the normalisation of the projection. In the frame and scaling of
Definition~\ref{def:h}, $g$ is $|N|^3|p||q||k|$ times Waleffe's geometric factor. A nonzero scaling
does not affect any statement about the vanishing of $C$.

\begin{theorem}[Closed form in the triad frame]
\label{thm:coeff}
For $p + q = k$, $N = p \times q$, and any real $s_k, s_p, s_q$,
\begin{equation}
  \label{eq:closed}
  g \;=\; -\,i\,(N\cdot N)^2\,\bigl(s_p|p| + s_q|q| + s_k|k|\bigr), \qquad
  C \;=\; \frac{i}{4}\,(N\cdot N)^2\,\bigl(s_p|p| + s_q|q| + s_k|k|\bigr)\bigl(s_p|p| - s_q|q|\bigr).
\end{equation}
\end{theorem}

\begin{proof}[Proof sketch]
Expand~\eqref{eq:Cdef} with Definition~\ref{def:h}. The terms quadratic in $i$ are each a triple
product $(N \times x)\cdot N = 0$, so $i$ survives only linearly and no use of $i^2 = -1$ is
needed. The chirality-free term is $(N \times p) \times (N \times q) = (N \cdot (p\times q))\,N$,
parallel to $N$, hence annihilated by pairing against $N \times k \perp N$; the remaining terms
collapse via $N \times (N \times x) = (N\cdot x)N - |N|^2 x$ with $N \perp p, q, k$. What is left
is a degree-$9$ polynomial identity in six integer variables and the three formal magnitudes,
which the proof assistant closes by \texttt{ring}.
\end{proof}

\begin{remark}[Independent verification, and the error it caught]
\label{rem:sign}
The prefactor and the phase of $C$ depend on the frame and scaling; the formal development proves
that in the common frame $N$ the coefficient is purely imaginary with no further hypothesis, and
that a global choice of frame (one $N(k)$ per wavevector) transforms differently under $k \to -k$.
The zero locus does not depend on either. Our first derivation of~\eqref{eq:closed} carried
$-s_k|k|$; a triad-by-triad comparison against a brute-force construction of $h^s$ and of $g$
from~\eqref{eq:Cdef} agreed in all eight chirality classes but with two magnitudes permuted ---
the signature of a sign inside a factor. The corrected form agrees with brute force to
$2.1\times10^{-16}$ in floating point and, because both sides are polynomials in $|p|,|q|,|k|$
with integer coefficients, is verified as an \emph{exact integer} polynomial identity over $5\,256$
triad-by-chirality instances, with an orientation flip $p \times q \to q \times p$ as the sharpest
of three negative controls shown to break it. A formal proof certifies the proof, not the
transcription of the statement; the exact-integer check certifies the statement.
\end{remark}

\begin{theorem}[Resonance is collinearity]
\label{thm:reso}
Let $p, q \in \R^3\setminus\{0\}$ with $k = p + q \ne 0$. Then $s_k|k| + s_p|p| + s_q|q| = 0$ for
some $(s_k, s_p, s_q) \in \{\pm1\}^3$ if and only if $p \times q = 0$.
\end{theorem}

\begin{proof}
The magnitudes are positive, so the signs are not all equal; up to a global sign, exactly one is
negative, and the equation is one of $|k| = |p| + |q|$, $|p| = |k| + |q|$, $|q| = |k| + |p|$ ---
the equality cases of the triangle inequality for $k = p + q$, $p = k + (-q)$, $q = k + (-p)$,
each of which holds iff the two summands are parallel and like-directed, i.e.\ iff $p \times q = 0$.
Conversely, if $p \times q = 0$ then $p$ and $q$ are parallel and one of the three equalities holds
with the signs chosen accordingly.
\end{proof}

The theorem uses nothing about the lattice; it is stated over $\R^3$ and applies to $\Z^3$ by
restriction. Its content for the decomposition is negative: the resonance condition selects no
triads beyond those already degenerate (no transverse plane), so it is not an independent locus.
On the lattice it was first established by exhaustive exact-integer enumeration over $558\,090$
triads before being proved.

\section{The vanishing locus, exactly}
\label{sec:closed}

\begin{theorem}[Vanishing locus of the helical coefficient]
\label{thm:converse}
For $p + q = k$ in $\Z^3$, $N = p \times q$, and real $s_k, s_p, s_q$: $C = 0$ if and only if
\[
  p \times q = 0 \quad\text{or}\quad s_p|p| + s_q|q| + s_k|k| = 0 \quad\text{or}\quad s_p|p| = s_q|q|.
\]
If moreover $p, q, k \ne 0$ and $(s_k, s_p, s_q) \in \{\pm1\}^3$, the middle alternative implies the
first, so
\[
  C = 0 \iff \text{the triad is collinear, or } s_p|p| = s_q|q| .
\]
\end{theorem}

\begin{proof}
By~\eqref{eq:closed}, $C$ is $\frac{i}{4}(N\cdot N)^2$ times two real factors; $\C$ has no zero
divisors, so $C = 0$ iff $N \cdot N = 0$ (i.e.\ $N = 0$, as $N \in \Z^3$) or one factor vanishes.
The second claim is Theorem~\ref{thm:reso}.
\end{proof}

The balance locus is genuinely present and non-degenerate: $p = (1,0,0)$, $q = (0,1,0)$ have
$|p| = |q|$ and $p \times q \ne 0$, and more generally every pair of modes on a common lattice
sphere with equal chirality lies on it. What Theorem~\ref{thm:converse} adds to the classical
sufficient conditions is the ``only if'': no further vanishing configuration exists. The three
sufficient conditions had first been proved separately, by degenerating the frame; the closed form
re-derives all three algebraically without mentioning a degenerate frame, and the formal
development checks on every build that the two routes agree --- a check the erroneous first
derivation of Remark~\ref{rem:sign} would have failed.

\section{The resonant-triad incidence structure on a periodic lattice}
\label{sec:torus}

We now leave the coefficient and consider only which triads exist. On a periodic lattice, every
ordered pair of nonzero modes $u, v$ with $u + v \ne 0$ closes a triad, and the incidence
structure is exactly solvable.

\begin{theorem}[Exact spectrum]
\label{thm:torus}
Let $G$ be a finite abelian group with no element of order $2$, $\Lambda = G \setminus \{0\}$,
$n := |\Lambda|$. For distinct $u, v \in \Lambda$ let $A(u,v)$ be the number of pairs (triad, ordered
pair of slots) in which $u$ and $v$ occupy two slots of a triad $a + b = c$ of elements of $\Lambda$
(so that the coincidence $v = 2u$ is counted, not excluded), and $A(u,u) := 0$. Then
\[
  A(u,v) \;=\; \begin{cases} 4, & u + v = 0,\\ 6, & \text{otherwise,}\end{cases}
  \qquad\text{i.e.}\qquad A \;=\; 6(J - I) - 2P,
\]
with $J$ the all-ones matrix and $P$ the involution $v \mapsto -v$ (fixed-point free by the
$2$-torsion hypothesis). Hence every row sum is $6n - 8$; a vector $f$ with $\sum_\Lambda f = 0$ and
$f(-v) = f(v)$ satisfies $Af = -8f$, and one with $f(-v) = -f(v)$ satisfies $Af = -4f$; the
spectrum of the degree-normalised matrix $A/(6n-8)$ is $\{1,\ -4/(6n-8),\ -8/(6n-8)\}$, and its
normalised gap $1 - \lambda_2$ tends to $1$ as $n \to \infty$.
\end{theorem}

\begin{proof}
The three slot-pair types contribute $2$, $2$, and $2 - 2[u + v = 0]$ respectively (the last
because the triads $(u, v, u+v)$ and $(v, u, u+v)$ require $u + v \ne 0$), giving the stated
weight; this is a finite case analysis that the proof assistant discharges. Row sums: $6(n-1)$
from $6(J-I)$, minus $2$ from the single entry of $P$, which lies off the diagonal since
$-u \ne u$. For the eigen-relations, $Jf = 0$ when $\sum f = 0$, so $Af = -6f - 2Pf$, and
$Pf = \pm f$. Normalising by the degree gives the spectrum; $\lambda_2 = -4/(6n-8) \to 0$.
\end{proof}

The hypothesis is necessary: with an element of order $2$, $u = -u$ and the $P$-term would fall on
the excluded diagonal, changing the row sum. The intended instance $(\Z/3)^3$ satisfies it; a
$2$-torsion counterexample is recorded in the formal development.

\begin{openproblem}[The spherical truncation]
\label{prob:56}
Define the same incidence weight on the spherical truncation $\Lam \subset \Z^3$, where a pair
$u, v$ closes a triad only if the third member also lies in the ball. The normalised gap
$1 - \lambda_2$ of the resulting matrix, computed exactly, is
\[
\begin{array}{c|cccccc}
M & 2 & 3 & 4 & 5 & 6 & 7\\ \hline
|\Lam| & 32 & 122 & 256 & 514 & 924 & 1418\\
1-\lambda_2 & 0.834985 & 0.833575 & 0.833419 & 0.833392 & 0.833359 & 0.833348\\
(1-\lambda_2) - 5/6 & 1.7\times10^{-3} & 2.4\times10^{-4} & 8.6\times10^{-5} & 5.9\times10^{-5} & 2.5\times10^{-5} & 1.5\times10^{-5}
\end{array}
\]
(six points, strictly decreasing; at $M = 7$ there are $939\,930$ triads). We conjecture
$1 - \lambda_2 \to 5/6$: the deviation from Theorem~\ref{thm:torus}'s limit $1$ would then be a
pure boundary effect of the sphere cutoff. Heuristically, in the continuum limit the matrix becomes
the kernel $2\cdot\mathbf 1_B(x+y) + 4\cdot\mathbf 1_B(x-y)$ on the unit ball $B$, and the
conjecture is that its second normalised eigenvalue is $1/6$; we have not proved this and would
welcome a direct attack.
\end{openproblem}

\section{Open problems and directions}
\label{sec:open}

\begin{enumerate}[leftmargin=*]
\item \textbf{The spherical gap (Open Problem~\ref{prob:56}).} An exactly solved periodic analogue
  sits next to a cleanly posed, numerically stable conjecture; the boundary correction may itself
  be exactly computable.
\item \textbf{The weighted identity as a tool.} Theorem~\ref{thm:weighted} uses only that the
  observable is a triad sum with one distinguished divergence-free index and a weight on the
  other two. The same computation should isolate the exact production mechanism, in the form
  of~\eqref{eq:enstrophy}, for other triad-structured quadratic systems --- shell models, surface
  quasi-geostrophic flow, other Galerkin truncations --- once the right weight is identified.
\item \textbf{Other bases.} Theorem~\ref{thm:converse} rests only on the factorisation
  \eqref{eq:closed} and the absence of zero divisors; the argument should transfer to other
  chirality-adapted bases and to $\R^3$, where Theorem~\ref{thm:reso} already holds.
\item \textbf{Bounds uniform in $M$.} Everything above is exact at fixed $M$ and gives no
  estimate. Whether the production term of~\eqref{eq:enstrophybalance} can be dominated by the
  dissipation uniformly in $M$ is the question that connects this system to global regularity for
  the untruncated equations, and it is open. One route in which exact identities of the present
  kind are the natural input is the self-consistent a-priori-bounds method of rigorous numerics
  for dissipative PDEs~\cite{zgliczynski2001}: assume an explicit tail envelope, show dissipation
  dominates the forcing of every tail mode, conclude invariance. We have not carried this out and
  make no claim about it.
\end{enumerate}

\appendix
\section*{Appendix: reproducibility}

The Lean~4 development, the exact-arithmetic cross-checks (Python, exact integers and rationals,
each negative control implemented and shown to fail), and the enumerations behind Open
Problem~\ref{prob:56} are in the public repository cited below.\footnote{\url{https://github.com/xaviercallens/SocrateAI-Scientific-MechanicaFluidorum},
commit \texttt{1aa60e0} or later.} The correspondence between statements and formal names is:
\begin{itemize}[leftmargin=*,itemsep=1pt]
\raggedright
\item Theorem~\ref{thm:weighted}: \texttt{AbstractAlgebraicConservation.lean},
  \texttt{weighted\_triad\_pairing} and \texttt{weighted\_triad\_sum}.
\item Corollaries~\ref{cor:energy}--\ref{cor:balance}: \texttt{FourierDynamicsZ3.lean},
  \texttt{energy\_conservation}, \texttt{enstrophy\_production\_identity},
  \texttt{energyRateZ3\_eq}, \texttt{enstrophyRateZ3\_eq}.
\item Theorems~\ref{thm:coeff}--\ref{thm:converse}: \texttt{HelicalBasis.lean},
  \texttt{gOf\_closed\_form}, \texttt{cOf\_closed\_form},
  \texttt{resonance\_implies\_collinear\_full}, \texttt{cOf\_eq\_zero\_iff}.
\item Theorem~\ref{thm:torus}: \texttt{TriadTorus.lean}, \texttt{A\_eq}, \texttt{row\_sum},
  \texttt{eigen\_even}, \texttt{eigen\_odd}.
\end{itemize}
Axiom footprints are reproduced by \texttt{\#print axioms} on each name.

\begin{thebibliography}{9}
\bibitem{kraichnan1959} R.~H. Kraichnan, The structure of isotropic turbulence at very high
Reynolds numbers, \emph{J. Fluid Mech.} \textbf{5} (1959) 497--543.
\bibitem{waleffe1992} F. Waleffe, The nature of triad interactions in homogeneous turbulence,
\emph{Phys. Fluids A} \textbf{4} (1992) 350--363.
\bibitem{zgliczynski2001} P. Zgliczy\'nski, K. Mischaikow, Rigorous numerics for partial
differential equations: the Kuramoto--Sivashinsky equation, \emph{Found. Comput. Math.}
\textbf{1} (2001) 255--288, DOI 10.1007/s002080010010.
\bibitem{lean4} L. de Moura, S. Ullrich, The Lean 4 theorem prover and programming language,
\emph{CADE-28} (2021).
\bibitem{mathlib} The mathlib Community, The Lean mathematical library, \emph{CPP 2020},
DOI 10.1145/3372885.3373824.
\end{thebibliography}

\end{document}
